%% tags: [KSAT - English, KSAT - English 2025]
%% type: mcq

\documentclass{article}
\usepackage[T1]{fontenc}
\usepackage{lmodern}
\usepackage[margin=1in]{geometry}

% no-op markers for the importer
\newenvironment{asset}[1]{}{}
\newcommand{\uses}[1]{}

\begin{document}

\begin{enumerate}

% 18
\item 다음 글의 목적으로 가장 적절한 것은? \answer{A}

Dear Rosydale City Marathon Racers,

We are really grateful to all of you who have signed up for the 10th Rosydale City Marathon that was scheduled for this coming Saturday at 10 a.m. Unfortunately, as you may already know, the weather forecast says that there is going to be a downpour throughout the race day. We truly hoped that the race would go smoothly. However, it is likely that the heavy rain will make the roads too slippery and dangerous for the racers to run safely. As a result, we have decided to cancel the race. We hope you understand and we promise to hold another race in the near future.

Sincerely,
Martha Kingsley, Race Manager

\begin{enumerate}
  \item 마라톤 경기 취소 사실을 공지하려고
  \item 마라톤 경기 사전 행사 참여를 독려하려고
  \item 마라톤 경기 참가비 환불 절차를 설명하려고
  \item 마라톤 경기 참여 시 규칙 준수를 당부하려고
  \item 마라톤 경기 진행에 따른 도로 통제를 안내하려고
\end{enumerate}

% 19
\item 다음 글에 드러난 Peter의 심경 변화로 가장 적절한 것은? \answer{B}

It was Valentine's Day on Friday and Peter was certain that his wife, Amy, was going to love his surprise. Peter had spent a long time searching online for an event that would be a new way to spend time with Amy. He had finally found the perfect thing for her. She often told him that she liked to go to places she had never visited before, and he was absolutely sure that she would love going to the new, five-star restaurant downtown. He smiled as he called the restaurant and asked for a reservation for Friday. Unfortunately, his smile quickly disappeared when he was told that the restaurant was fully reserved. "That's too bad," he said quietly. "I thought that I had found the right place."

\begin{enumerate}
  \item relaxed $\to$ indifferent
  \item confident $\to$ disappointed
  \item confused $\to$ satisfied
  \item jealous $\to$ discouraged
  \item embarrassed $\to$ joyful
\end{enumerate}

% 20
\item 다음 글에서 필자가 주장하는 바로 가장 적절한 것은? \answer{A}

We almost universally accept that playing video games is at best a pleasant break from a student’s learning and more often what prevents a student from accomplishing their goals. Games catch and hold attention in a way that few things can. And yet once they have our focus, they rarely seem to offer anything meaningful to help students grow in their lives outside the games. While this may be true for many games, we are too easily ignoring a valuable tool that could be used to enhance productivity instead of derailing it. Rather, it is desirable that we develop games that connect to the learning outcomes we want for our students. This will enable educators to take advantage of games’ attention commanding capacities and allow our students to enjoy their games while learning.

\begin{enumerate}
  \item 학습 효과 증진에 활용될 수 있는 게임을 개발해야 한다.
  \item 교육 현장에서 학습과 게임 활동을 적절하게 분배해야 한다.
  \item 학습 활동에 게임이 초래하는 집중력 저하를 경계해야 한다.
  \item 여가 시간에 게임을 활용함으로써 학습 효율을 향상해야 한다.
  \item 게임의 부정적 영향을 줄이기 위해 학습 공동체가 노력해야 한다.
\end{enumerate}

% 21
\item 밑줄 친 "hunting the shadow, not the substance"가 의미하는 바로 가장 적절한 것은? \answer{C}

The passage explains Vitruvius’ view that architecture requires both practical and theoretical knowledge; relying only on theories and scholarship is like "hunting the shadow, not the substance."

\begin{enumerate}
  \item seeking abstract knowledge emphasized by architectural tradition
  \item discounting the subjects necessary to achieve architectural goals
  \item pursuing the ideals of architecture without the practical skills
  \item prioritizing architecture’s material aspects over its artistic ones
  \item following historical precedents without regard to current standards
\end{enumerate}

% 22
\item 다음 글의 요지로 가장 적절한 것은? \answer{B}

The ability to understand emotions — to have a diverse emotion vocabulary and to understand the causes and consequences of emotion — is particularly relevant in group settings. Individuals skilled in this domain can express emotions accurately, facilitating clear communication and cooperation, appreciate others’ perspectives, and promote constructive conflict resolution and improved group functioning.

\begin{enumerate}
  \item 집단 구성원 간 갈등 해소를 위해 감정 조절이 중요하다.
  \item 감정 이해 능력은 집단 내 원활한 소통과 협력을 촉진한다.
  \item 타인에 대한 공감 능력은 자신의 감정 표현 능력을 향상한다.
  \item 감정 관련 어휘에 대한 지식은 공감 능력 발달의 기반이 된다.
  \item 자신의 감정 상태에 대한 이해는 사회성 함양에 필수적 요소이다.
\end{enumerate}

% 23
\item 다음 글의 주제로 가장 적절한 것은? \answer{A}

The Industrial Age changed the relationship among time, labor, and capital: around-the-clock factory shifts tied wages to effort and production and reorganized society around new principles of productivity.

\begin{enumerate}
  \item shift in the work-time paradigm brought about by industrialization
  \item effects of standardizing production procedures on labor markets
  \item influence of industrialization on the machine-human relationship
  \item efficient ways to increase the value of time in the Industrial Age
  \item problems that excessive work hours have caused for laborers
\end{enumerate}

% 24
\item 다음 글의 제목으로 가장 적절한 것은? \answer{E}

The selfie extends the history of the self-portrait, becoming a key visual signature of the networked era by enabling more people to depict themselves and express the drama of daily self-presentation.

\begin{enumerate}
  \item Are Selfies Just a Temporary Trend in Art History?
  \item Fantasy or Reality: Your Selfie Is Not the Real You
  \item The Selfie: A Symbol of Self-oriented Global Culture
  \item The End of Self-portraits: How Selfies Are Taking Over
  \item Selfies, the Latest Innovation in Representing Ourselves
\end{enumerate}

% 25
\item 다음 도표의 내용과 일치하지 않는 것은? \answer{A} % adjust key if needed after real chart

\begin{enumerate}
  \item For each year, producers had the highest percentage among the four roles.
  \item Directors’ percentage in 2021 was lower than 2020 but higher than 2022.
  \item Writers increased by 4pp from 2020 to 2021 and by 1pp from 2021 to 2022.
  \item Editors were less than 20% in each of the three years.
  \item In 2022, producers’ percentage was the same as in 2020.
\end{enumerate}

% 26
\item Dick Enberg에 관한 다음 글의 내용과 일치하지 않는 것은? \answer{A} % placeholder key

\begin{enumerate}
  \item Michigan에서 태어났다.
  \item 대학 야구팀 코치였다.
  \item 중국을 방문한 첫 미국인 스포츠 캐스터였다.
  \item 마지막 생방송 후 3년 뒤에 사망하였다.
  \item Emmy Awards를 수상하였다.
\end{enumerate}

% 27
\item Adenville City Pass Card에 관한 다음 안내문의 내용과 일치하지 않는 것은? \answer{C}

\begin{enumerate}
  \item 관광객을 위한 대중교통 카드이다.
  \item 시티 투어 버스에는 사용할 수 없다.
  \item 5일 패스 카드에만 주요 관광지 입장료 할인 혜택이 제공된다.
  \item 미사용 카드는 구입일로부터 30일 이내에 환불이 가능하다.
  \item 모바일 카드는 A-Transit 앱에서 구입할 수 있다.
\end{enumerate}

% 28
\item Luckwood Snow Festival에 관한 다음 안내문의 내용과 일치하는 것은? \answer{C}

Festival Info  
15th annual Luckwood Snow Festival. Dates: Jan 24–30 (7 days), 9 a.m.–8 p.m. Venue: Luckwood Park.  
Special Activities: Snow Sculpture Contest (11 teams); Kids’ tunnels and slides.  
Transportation: No parking (use public transit/shuttle). Shuttle runs between Luckwood Subway Station and Luckwood Park (one-way fare: $1, cash only).

\begin{enumerate}
  \item 2년에 한 번 열린다.
  \item 열흘 동안 진행된다.
  \item 눈 조각 경연에는 11개 팀이 참가할 것이다.
  \item 주차가 가능하다.
  \item 셔틀버스 이용은 무료이다.
\end{enumerate}

% 29
\item \textit{다음 글의 밑줄 친 부분 중, 어법상 틀린 것은?} \answer{B}

Think of yourself. When you decide to get up and get a drink of water, for example, you don’t consciously organize or consider the host of steps involved. Imagine if we \underline{had to} consider every single muscle that needed to be contracted or relaxed just to stand up and walk. It would be tiresome and very slow — as patients recovering from a brain injury affecting the motor system \underline{knows}. The autopilot parts of our brain do it for us automatically, \underline{freeing up} our conscious mind for more important jobs. It is the older parts of our brain \underline{that support these automatic processes that allow} us to move, hear, see, and use many of our social skills. More recently evolved abilities like talking, reading, and writing are far less automated. So, most of the time, \underline{what you are perceiving, feeling, or thinking is} based on a very crude and fast analysis that happens completely without your awareness.

\begin{enumerate}
  \item ① had to
  \item ② knows
  \item ③ freeing up
  \item ④ that support these automatic processes that allow
  \item ⑤ what you are perceiving, feeling, or thinking is
\end{enumerate}

% 30
\item \textit{다음 글의 밑줄 친 부분 중, 문맥상 낱말의 쓰임이 적절하지 않은 것은?} \answer{D}

Studies in psychology have reported cases in which competitive incentives resulted in lower task effort, and their focus was on the psychological underpinnings of the reduction in motivation. For example, competition presents an inevitable conflict between the motivation to achieve one’s personal goal and the \underline{desire} to maintain good relationships with others. When the maintenance of interpersonal relationships is important, competitors experience an \underline{internal} conflict that can harm their desire to achieve their goal and taint the good feeling brought about by winning. Exline and Lobel found that the perception of oneself as a target for upward social comparison often makes people \underline{uncomfortable}. When they believe that others are making envious comparisons with them, people feel uneasiness, distress, or sorrow. Feelings of guilt, an emotion generally associated with high motivation for goal-achievement, lead to \underline{stronger} motivation and performance in the pursuit of competitive goals. Consequences of this emotional state include lower task motivation in a competition and preferences for more cooperative and altruistic outcomes, such as \underline{diminishing} the significance of the outcome or sharing the winner’s reward.

\begin{enumerate}
  \item ① desire
  \item ② internal
  \item ③ uncomfortable
  \item ④ stronger
  \item ⑤ diminishing
\end{enumerate}

% 31
\item \textit{다음 빈칸에 들어갈 말로 가장 적절한 것을 고르시오.} \answer{E}

Literature can be helpful in the language learning process because of the it fosters in readers. Core language teaching materials must concentrate on how a language operates both as a rule-based system and as a sociosemantic system. Very often, the process of learning is essentially analytic, piecemeal, and, at the level of the personality, fairly superficial. Engaging imaginatively with literature enables learners to shift the focus of their attention beyond the more mechanical aspects of the foreign language system. When a novel, play or short story is explored over a period of time, the result is that the reader begins to ‘inhabit’ the text. He or she is drawn into the book. Pinpointing what individual words or phrases may mean becomes less important than pursuing the development of the story. The reader is eager to find out what happens as events unfold; he or she feels close to certain characters and shares their emotional responses. The language becomes ‘transparent’ — the fiction draws the whole person into its own world.

\begin{enumerate}
  \item linguistic insight
  \item artistic imagination
  \item literary sensibility
  \item alternative perspective
  \item personal involvement
\end{enumerate}

% 32
\item \textit{다음 빈칸에 들어갈 말로 가장 적절한 것을 고르시오.} \answer{C}

Education, at its best, teaches more than just knowledge. It teaches critical thinking: the ability to stop and think before acting, to avoid succumbing to emotional pressures. This is not thought control. It is the very reverse: mental liberation. Even the most advanced intellectual will be imperfect at this skill. But even imperfect possession of it \underline{\hspace{1.5em}} of being ‘stimulus-driven’, constantly reacting to the immediate environment, the brightest colours or loudest sounds. Being driven by heuristic responses, living by instinct and emotion all the time, is a very easy way to live, in many ways: thought is effortful, especially for the inexperienced. But emotions are also exhausting, and short-term reactions may not, in the long term, be the most beneficial for health and survival. Just as we reach for burgers for the sake of convenience, storing up the arterial fat which may one day kill us, so our reliance on feelings can do us great harm.

\begin{enumerate}
  \item intensifies people’s danger
  \item enhances our understanding
  \item frees a person from the burden
  \item allows us to accept the inevitability
  \item requires one to have the experience
\end{enumerate}

% 33
\item \textit{다음 빈칸에 들어갈 말로 가장 적절한 것을 고르시오.} \answer{B}

We are famously living in the era of the attention economy, where the largest and most profitable businesses in the world are those that consume my attention. The advertising industry is literally dedicated to capturing the conscious hours of my life and selling them to someone else. It might seem magical that so many exciting and useful software systems are available to use for free, but it is now conventional wisdom that if you can’t see who is paying for something that appears to be free, then \underline{\hspace{5em}}.

\begin{enumerate}
  \item all of your attention has already been spent
  \item the real product being sold is you
  \item your privacy is being violated
  \item the public may be sponsoring you
  \item you owe the benefits to your friend AI
\end{enumerate}

% 34
\item \textit{다음 빈칸에 들어갈 말로 가장 적절한 것을 고르시오.} \answer{E}

Centralized, formal rules can \underline{\hspace{6em}}. The rules of baseball don’t just regulate the behavior of the players; they determine the behavior that constitutes playing the game. … Without them an individual would not have the opportunity to occupy the role.

\begin{enumerate}
  \item categorize one’s patterns of conduct in legal and productive ways
  \item lead people to reevaluate their roles and practices in a society
  \item encourage new ways of thinking which promote creative ideas
  \item reinforce one’s behavior within legal and established contexts
  \item facilitate productive activity by establishing roles and practices
\end{enumerate}

% 35
\item \textit{다음 글에서 전체 흐름과 관계 없는 문장은?} \answer{D}

(The expansion of sports tourism …)  
\begin{enumerate}
  \item ① The significance of the car … has had no less an impact on sports tourism specifically.
  \item ② 1920s in the USA … later in Britain.
  \item ③ access to many areas not served by public transport …
  \item ④ The expansion of reasonably priced, good quality accommodation … restaurants.
  \item ⑤ it was invaluable for the development of many forms of sports tourism …
\end{enumerate}

% 36
\item \textit{주어진 글 다음에 이어질 글의 순서로 가장 적절한 것을 고르시오.} \answer{A}

(지문 요약: 농지 임대에서 평판이 시장 집행을 가능케 함)

\begin{enumerate}
  \item (C) → (B) → (A)
  \item (B) → (C) → (A)
  \item (A) → (B) → (C)
  \item (B) → (A) → (C)
  \item (C) → (A) → (B)
\end{enumerate}

% 37
\item \textit{주어진 글 다음에 이어질 글의 순서로 가장 적절한 것을 고르시오.} \answer{A}

(지문 요약: 원형 배열 새들이 선형 배열보다 감정 전염/스캐닝이 더 조정됨)

\begin{enumerate}
  \item (B) → (C) → (A)
  \item (C) → (B) → (A)
  \item (A) → (B) → (C)
  \item (A) → (C) → (B)
  \item (C) → (A) → (B)
\end{enumerate}

% 38
\item \textit{글의 흐름으로 보아, 주어진 문장이 들어가기에 가장 적절한 곳을 고르시오.} \answer{E}

주어진 문장: \emph{“Without any special legal protection for trade secrets, however, the secretive inventor risks that an employee or contractor will disclose the proprietary information.”}

\begin{enumerate}
  \item ①
  \item ②
  \item ③
  \item ④
  \item ⑤
\end{enumerate}

% 39
\item \textit{글의 흐름으로 보아, 주어진 문장이 들어가기에 가장 적절한 곳을 고르시오.} \answer{D}

주어진 문장: \emph{“In reality, objects do not conform to a linear lifecycle model; instead, they undergo breakdowns, await repairs, are stored away, or are repurposed later.”}

\begin{enumerate}
  \item ①
  \item ②
  \item ③
  \item ④
  \item ⑤
\end{enumerate}

% 40
\item \textit{다음 글을 한 문장으로 요약하고, 빈칸 (A), (B)에 들어갈 말로 가장 적절한 것은?} \answer{A}

\begin{enumerate}
  \item (A) controllability \quad (B) challenge
  \item (A) predictability \quad (B) support
  \item (A) manageability \quad (B) reverse
  \item (A) affordability \quad (B) intensify
  \item (A) accessibility \quad (B) question
\end{enumerate}

% 41
\item \textit{윗글의 제목으로 가장 적절한 것은?} \answer{B}

\begin{enumerate}
  \item Anatomical Distance Between Humans and Other Primates
  \item Human Hands: A Decisive Leap in the Evolutionary Path
  \item Our Hands: An Unexpected Outcome of Evolution
  \item Human Grip: The Dilemma of Human Survival
  \item Hidden Power of the Daily Use of Tools
\end{enumerate}

% 42
\item \textit{밑줄 친 (a)～(e) 중에서 문맥상 낱말의 쓰임이 적절하지 않은 것은?} \answer{E}

\begin{enumerate}
  \item (a)
  \item (b)
  \item (c)
  \item (d)
  \item (e)
\end{enumerate}

% 43
\item \textit{주어진 글 (A)에 이어질 내용을 순서에 맞게 배열한 것으로 가장 적절한 것은?} \answer{E}

\begin{enumerate}
  \item (B) → (C) → (D)
  \item (D) → (C) → (B)
  \item (C) → (B) → (D)
  \item (D) → (B) → (C)
  \item (C) → (D) → (B)
\end{enumerate}

% 44
\item \textit{밑줄 친 (a)～(e) 중에서 가리키는 대상이 나머지 넷과 다른 것은?} \answer{C}

\begin{enumerate}
  \item (a)
  \item (b)
  \item (c)
  \item (d)
  \item (e)
\end{enumerate}

% 45
\item \textit{윗글에 관한 내용으로 적절하지 않은 것은?} \answer{D}

\begin{enumerate}
  \item Grace는 Ethan과 Sean의 관계를 부러워했다고 말했다.
  \item Grace는 Ethan에게 Sean과 둘이서 하이킹할 것을 권했다.
  \item Sean의 선글라스가 서랍장 안에 있었다.
  \item Ethan은 혼자서 세탁소에 하이킹 재킷을 찾으러 갔다.
  \item Sean은 하이킹하는 도중 돌에 걸려 넘어졌다.
\end{enumerate}


\end{enumerate}

\end{document}



    


